\section{Interfacial tails and accuracy diagnostics}
\label{sec:capillarity-diagnostics}

This appendix gives the capillarity model behind the large-$\gamma$
range~\eqref{eq:two-dimensional-boundary-range} of the
two-dimensional boundary law and the assumption $\alpha>1/2$ in the
density-envelope route to the boundary law ($\alpha=1/2$ is the
two-dimensional interfacial exponent). It then contrasts full-state TV
with histogram norms and barrier accuracy.

\subsection{A capillarity model for the two-dimensional boundary}

Biskup, Chayes, and Koteck\'y established the competition between
background fluctuations and droplet surface
costs~\cite{BiskupChayesKotecky2002}. Leung and Zia located the
droplet-to-strip transition on a periodic square at fixed excess
volume fraction~\cite{LeungZia1990}, Neuhaus and Hager analyzed it for
first-order simulations~\cite{NeuhausHager2003}, and Kim, Keyes, and
Straub exhibited droplet and strip energy regions in a toroidal
two-dimensional Potts system~\cite{KimKeyesStraub2011}. These
results motivate the model below; none supplies our assumed full
density envelope.

Assume $Y_N=(E-e_{-,N})/\Delta_N$ is supported in a fixed compact
interval $K$ containing $[0,1]$ and satisfies a large-deviation
principle with speed $\sqrt N$ and good rate $I$; the model is stated
for $d=2$, where this speed is both the surface scale $N^{(d-1)/d}$ and
the fluctuation scale.
Assume also that $I$ is continuous and finite on $[0,1]$, infinite on
$K\setminus[0,1]$, zero exactly at $0,1$, and positive on $(0,1)$.
Compact support is a simplification, not a claim about systems with
unbounded kinetic energy.

\begin{proposition}[Capillarity variational law]
\label{prop:capillarity-ldp}
For $\Delta_N/N\to\ell>0$ and $c_N/N^{3/2}\to\gamma>0$, reweighting
by the exact reservoir~\eqref{eq:bath-dos} with the secant
calibration~\eqref{eq:secant-calibration} gives a large-deviation
principle at speed $\sqrt N$ with rate
\begin{equation}
 I_\gamma(\theta)=I(\theta)-J_\gamma(\theta)
 -\min_{0\le u\le1}\{I(u)-J_\gamma(u)\},
 \qquad
 J_\gamma(\theta)=\frac{\beta^2\ell^2}{2\gamma}\theta(1-\theta).
 \label{eq:capillarity-rate}
\end{equation}
If
\begin{equation}
 \gamma_*=\frac{\beta^2\ell^2}{2}
 \sup_{0<\theta<1}\frac{\theta(1-\theta)}{I(\theta)}<\infty,
 \label{eq:capillarity-threshold}
\end{equation}
then $\gamma>\gamma_*$ leaves only the endpoints as rate minimizers.
For $\gamma<\gamma_*$ all minimizing energies lie away from the
endpoints and $\TV(P_N,Q_N)\to1$.
\end{proposition}
\begin{proof}
Equation~\eqref{eq:physical-interior-exact}, expanded on $K$, gives
$h_N/\sqrt N\to J_\gamma$ uniformly; the cutoff lies beyond $K$ for
large $N$ because $c_N/\Delta_N\to\infty$. For this continuous
bounded tilt, the large-deviation upper bound on a finite cover by
sets where the tilt varies by at most $\epsilon$, with the lower
bound on an open neighborhood of each point, gives
\[
 \lim_N\frac1{\sqrt N}\log\E e^{h_N}
 =\max_{\theta\in K}\{J_\gamma(\theta)-I(\theta)\}.
\]
The same bounds for restricted integrals prove
\eqref{eq:capillarity-rate}. If $\gamma>\gamma_*$,
$I(\theta)>J_\gamma(\theta)$ throughout the interior.
If $\gamma<\gamma_*$, the minimum of $I-J_\gamma$ is strictly
negative, whereas its endpoint values are zero, so by compactness
and continuity its minimizing set lies a positive distance from the
endpoints. A neighborhood of that set has reservoir probability
tending to one and canonical probability tending to zero, so
$\TV(P_N,Q_N)\to1$.
\end{proof}

For $\gamma>\gamma_*$ the proposition gives macroscopic energy
concentration, not full-state TV or the fluctuation law within a
phase. At $\gamma=\gamma_*$ the model does not identify the weights
of tied minimizers: multiplying the law near them by positive
constants, or by factors with logarithm $o(\sqrt N)$, leaves the rate
unchanged but can change those weights or select one minimizer.

For an explicit square-torus model, set
\begin{equation}
 I(\theta)=\tau\min\{2\sqrt{\pi\theta},\,2,\,
                         2\sqrt{\pi(1-\theta)}\},
 \qquad 0\le\theta\le1,\quad \tau>0.
 \label{eq:square-torus-model}
\end{equation}
The branches are an isotropic circular droplet, a slab with two
interfaces, and the complementary bubble: the isotropic-tension form
of the standard square-torus construction
\cite{LeungZia1990,NeuhausHager2003}, with the droplet-to-strip
crossover at $\theta=1/\pi$. Equation~\eqref{eq:square-torus-model}
is a model input, not a Wulff law derived for the lattice Potts
model.

\begin{corollary}[Square-torus minimizers]
\label{cor:square-capillarity}
For \eqref{eq:square-torus-model},
$\gamma_*=\beta^2\ell^2/(16\tau)$.
Above $\gamma_*$ the only minimizers are $0,1$; below it the unique
minimizer is $1/2$; at equality the minimizers are $0,1/2,1$.
\end{corollary}
\begin{proof}
For all $\theta$,
\begin{equation}
 I(\theta)\ge8\tau\theta(1-\theta),
 \label{eq:square-perimeter-bound}
\end{equation}
with equality exactly at $0,1/2,1$.
The slab branch obeys this because $\theta(1-\theta)\le1/4$.
For the droplet branch, away from $\theta=0$ the bound follows
from
$8\sqrt\theta(1-\theta)\le16/(3\sqrt3)<2\sqrt\pi$;
the bubble branch is symmetric. Equality at $1/2$ gives the stated
$\gamma_*$.
Put $A=\beta^2\ell^2/(2\gamma)$.
If $A<8\tau$, only endpoints minimize $I-A\theta(1-\theta)$.
If $A=8\tau$, equality in \eqref{eq:square-perimeter-bound} identifies
the three minimizers. If $A>8\tau$, using $I(1/2)=2\tau$ gives
\[
 [I(\theta)-A\theta(1-\theta)]-[I(1/2)-A/4]
 \ge(A-8\tau)(\theta-1/2)^2,
\]
which is strictly positive off $1/2$.
\end{proof}

\subsection{Histogram norms can hide a persistent error}

An unscaled density norm can vanish while TV stays fixed, even within
one phase. Let $p_\sigma$ be the density of $\Normal(0,\sigma^2)$
and $q_\sigma$ that of $\Normal(0,r\sigma^2)$, for fixed $0<r<1$.
Gaussian integration gives
\begin{equation}
 \|p_\sigma-q_\sigma\|_2^2
 =\frac1{2\sqrt\pi\,\sigma}
 \left[1+r^{-1/2}-2\sqrt{\frac2{1+r}}\right]\longrightarrow0
 \quad(\sigma\to\infty),
 \label{eq:histogram-l2}
\end{equation}
although
\begin{equation}
 \TV(p_\sigma,q_\sigma)
 =2\left[
 \Phi\!\left(\sqrt{\frac{\log r}{r-1}}\right)
 -\Phi\!\left(\sqrt{\frac{r\log r}{r-1}}\right)\right]>0.
 \label{eq:histogram-tv}
\end{equation}
The densities cross at $|E|/\sigma=\sqrt{r\log r/(r-1)}$, and
integrating their difference between the crossings
gives~\eqref{eq:histogram-tv}. For histogram probabilities in bins of
fixed width $a>0$, a rescaled Riemann sum gives
$\sum_j(P_j-Q_j)^2\sim a\|p_\sigma-q_\sigma\|_2^2$, so the Euclidean
histogram norm decays as $\sigma^{-1/2}$, or $N^{-1/4}$ when
$\sigma\asymp\sqrt N$, although the variance ratio and the TV error
stay fixed.

This is the distinction promised in Section~\ref{sec:introduction}:
the finite-bath Potts comparison of Griffin, Matty, and Swendsen
measures agreement by the unweighted Euclidean norm of
energy-probability differences, their equation~(23), and reports its
decay with size \cite{GriffinMattySwendsen2017}; the example shows
that such decay is compatible with a fixed full-state error.

\subsection{Barrier accuracy requires a separate scale}

Let the canonical and reservoir energy densities $p$ and $q$ be
positive at specified energies $e_-$ and $E^\dagger$, with
fixed-energy log barriers $B_P=\log[p(e_-)/p(E^\dagger)]$ and
$B_Q=\log[q(e_-)/q(E^\dagger)]$.
With $h(e_-)=0$, normalization cancels exactly:
\begin{equation}
 B_Q-B_P=-h(E^\dagger).
 \label{eq:fixed-barrier-shift}
\end{equation}
The same holds for ratios of point probabilities in a discrete
energy law. For a quadratic weight balanced at the endpoints with gap
$\Delta=2m$, the midpoint gain is $\kappa m^2/2=\kappa\Delta^2/8$; for
the exact secant-calibrated reservoir,
\begin{equation}
 h\!\left(\frac{e_-+e_+}{2}\right)
 =c\log\cosh\!\left(\frac{\beta\Delta}{2c}\right),
 \label{eq:physical-exact-barrier}
\end{equation}
which follows from \eqref{eq:physical-interior-exact} at
$\theta=1/2$ and is positive.

If $c\gg\Delta$, the barrier lowering is
$\beta^2\Delta^2/(8c)\,[1+O((\Delta/c)^2)]$.
For $\Delta\to\infty$ and inverse temperature bounded
away from zero and infinity, the lowering vanishes if and only if
$c\gg\Delta^2$, because $\log\cosh x$ lies between a positive
constant times $\min(x^2,|x|)$ and $x^2/2$; the lower bound follows
from continuity of the corresponding ratios on $0<|x|\le1$ and
$|x|\ge1$, with positive limits at zero and infinity. Along a
subsequence with $\Delta/c$ bounded below the barrier error is thus
of order at least $\Delta$, while in the other regime it is
comparable to $\Delta^2/c$.

For an extensive gap, absolute accuracy of the midpoint barrier
therefore requires $c\gg N^2$, whereas the two-phase full-state TV
error can vanish for $c\gg N^{3/2}$. Relative accuracy is less
demanding: for a canonical barrier $B_P\asymp N^{(d-1)/d}$ and
$c\gg N$, it requires $c\gg N^2/B_P\asymp N^{1+1/d}$, at most
$N^{3/2}$ for $d\ge2$. These statements concern density ratios at
specified energies, not shifted peaks, saddle points, or dynamical
transition rates.
